\mychapter{0}{Introducción}
\label{Introducción}

Los robots humanoides han pasado en pocos años de ser plataformas de investigación costosas y
frágiles a productos comerciales con un precio comparable al de un vehículo industrial. En paralelo,
el aprendizaje de políticas de manipulación a partir de demostraciones humanas ha desplazado en
buena medida a la programación explícita de trayectorias. La confluencia de ambas líneas plantea una
pregunta práctica que este trabajo aborda de forma acotada y medible: cuando se dispone de un
conjunto de demostraciones de una tarea industrial concreta, ¿compensa emplear un modelo fundacional
de gran tamaño, o basta con una arquitectura clásica de imitación entrenada desde cero?

Este Trabajo de Fin de Máster responde a esa pregunta sobre una tarea única y bien delimitada —la
apertura de una válvula de volante mediante el humanoide Unitree G1, íntegramente en simulación—
comparando dos métodos entrenados sobre exactamente el mismo conjunto de datos y evaluados en
exactamente las mismas condiciones.


\section*{Planteamiento del problema}

En las instalaciones del sector del petróleo y el gas, la operación manual de válvulas de volante es
una tarea rutinaria, físicamente exigente y a menudo situada en zonas clasificadas como atmósferas
potencialmente explosivas. Es, por su naturaleza repetitiva y por el riesgo asociado al entorno, una
de las candidatas naturales a la automatización mediante robots móviles de propósito general. Sin
embargo, se trata de una tarea que resiste mal el enfoque clásico de programar la trayectoria: la
válvula puede estar montada en cualquier orientación, el volante exige un agarre firme y un giro
sostenido, y la reacción del mecanismo se propaga por los brazos hasta comprometer el equilibrio del
robot.

El aprendizaje por imitación ofrece una alternativa. En lugar de describir cómo debe moverse el
robot, se le muestra la tarea un cierto número de veces y se ajusta una política que reproduzca ese
comportamiento. El coste se traslada entonces de la ingeniería de la trayectoria a la recogida de
demostraciones, que en el caso de un humanoide de veintinueve grados de libertad no es trivial.

Dentro de ese marco conviven hoy dos aproximaciones de naturaleza muy distinta. Por un lado, los
modelos visión-lenguaje-acción (VLA), que parten de un modelo fundacional preentrenado sobre
grandes volúmenes de datos heterogéneos y se ajustan a la tarea concreta; NVIDIA GR00T
\cite{nvidia2025gr00t} es un exponente reciente orientado específicamente a humanoides. Por otro,
las arquitecturas de imitación específicas, entrenadas desde cero sobre los datos de la tarea, de
las que ACT \cite{zhao2023aloha} es probablemente la más extendida. La diferencia de escala entre
ambas es de casi dos órdenes de magnitud en número de parámetros.

La literatura ofrece abundantes resultados de cada familia por separado, pero comparaciones
rigurosas entre ambas sobre una misma tarea, un mismo conjunto de datos y un mismo criterio de
éxito son escasas. Y sin esa comparación, la decisión de qué método emplear en un despliegue
concreto se toma sobre intuiciones acerca del tamaño del modelo, no sobre evidencia.


\section*{Objetivos}

El objetivo general del trabajo es \textbf{cuantificar, sobre una tarea de manipulación industrial
representativa y en igualdad de condiciones, la diferencia de rendimiento entre un modelo
visión-lenguaje-acción y una política de imitación clásica}.

De él se derivan los siguientes objetivos específicos:

\begin{enumerate}

	\item Construir un entorno de simulación reproducible de la tarea, con el humanoide Unitree G1
	y una válvula procedente de un modelo CAD real, que sirva a la vez para grabar demostraciones y
	para evaluar políticas con un criterio de éxito objetivo.

	\item Recoger un conjunto de demostraciones teleoperadas con variabilidad suficiente para que
	la tarea no sea memorizable, incluyendo dos disposiciones distintas de la válvula y
	aleatorización de su posición.

	\item Establecer una única cadena de conversión que alimente a ambos métodos desde el mismo
	fichero de demostraciones, de modo que ninguna diferencia observada pueda atribuirse al
	tratamiento de los datos.

	\item Entrenar y evaluar las dos políticas, y contrastar sus tasas de éxito con la prueba
	estadística adecuada a un diseño de medidas emparejadas.

	\item Caracterizar el coste de cada método en los dos ejes que condicionan un despliegue real:
	la cantidad de demostraciones necesaria y la latencia de inferencia frente al presupuesto de
	tiempo del lazo de control.

\end{enumerate}


\section*{Metodología}

El trabajo sigue una metodología experimental comparativa, articulada en cuatro etapas encadenadas:
teleoperación y grabación, conversión del conjunto de datos, entrenamiento y evaluación. El diseño
descansa sobre tres decisiones que condicionan la validez de las conclusiones y que conviene
adelantar aquí.

\textbf{Un único conjunto de datos para ambos métodos.} El formato de conjunto de datos de GR00T es
un superconjunto del formato LeRobot que consume ACT, de modo que una sola conversión sirve a los
dos entrenamientos. Cualquier diferencia de resultado procede, por tanto, del método y no de una
preparación distinta de los datos.

\textbf{Un único banco de evaluación.} Ambas políticas se sirven como procesos independientes que
implementan el mismo protocolo de comunicación con el simulador. El comando que lanza la evaluación
es idéntico hasta el último carácter para las dos; lo único que cambia es qué proceso atiende la
petición. Esto elimina toda una clase de diferencias espurias derivadas del montaje experimental.

\textbf{Comparación emparejada.} Las dos políticas se evalúan sobre las mismas condiciones iniciales,
generadas con la misma semilla, de manera que cada episodio constituye un par. El contraste
estadístico apropiado es entonces la prueba exacta de McNemar sobre los casos discordantes, no una
comparación de proporciones independientes, que desperdiciaría la información del emparejamiento y
sobreestimaría la incertidumbre.

Sobre esa base se realizan tres experimentos: la comparación principal sobre doscientas tiradas por
política, una curva de eficiencia en datos con subconjuntos estratificados de veinticinco, cincuenta
y cien demostraciones, y una medida de latencia de inferencia en lazo cerrado.

Todo el desarrollo se ha realizado en simulación. No se ha dispuesto de acceso a un robot físico, lo
que constituye la principal limitación del trabajo y se discute explícitamente en el capítulo de
evaluación.


\section*{Estructura del trabajo}

El resto de la memoria se organiza como sigue.

El capítulo \ref{Estudio del problema}, \textit{Estudio del problema}, sitúa el trabajo en su
contexto industrial y tecnológico, revisa el estado de la cuestión en aprendizaje por imitación,
modelos visión-lenguaje-acción, simulación robótica y teleoperación, delimita el problema concreto
que se aborda y recoge la normativa aplicable a un sistema de estas características.

El capítulo \ref{Gestión de proyecto software}, \textit{Gestión de proyecto software}, presenta la
simulación de la gestión del proyecto en condiciones de entorno empresarial: alcance, estimación,
presupuesto, planificación, recursos y riesgos.

El capítulo \ref{Solución}, \textit{Solución}, describe el sistema construido: la arquitectura, el
diseño de la tarea de simulación, la cadena de datos, los dos entrenamientos y las pruebas
realizadas sobre cada eslabón.

El capítulo \ref{Evaluación}, \textit{Evaluación}, expone el protocolo experimental y analiza los
resultados obtenidos, con especial atención a qué afirmaciones sostienen los datos y cuáles no.

El capítulo de \textit{Conclusión} recoge las aportaciones del trabajo, los problemas encontrados y
las líneas de continuación.

Cierran el documento tres anexos: el control de versiones del proyecto, el seguimiento real del
Trabajo de Fin de Máster y el protocolo de evaluación empleado.
